\addchap{Nomenclature}
\label{chap:nomenclature}

% Selection: recurring quantities with at least three uses in the active
% main manuscript. Appendices, comments and inactive sources are excluded.
% Indexed variants share a base entry; formulation-specific parameters remain
% defined locally. The notation key also covers recurring figure notation.
The following tables collect recurring quantities and notation. Reference
states, heat-flux signs and formulation-specific parameters are defined where
used. A dash denotes a dimensionless quantity. Symbols with more than one
meaning are distinguished by their context and units.

\begingroup
% The inside-page geometry is already active. Avoid resetting it from the
% TUM page-style hook while longtable is shipping a continuation page.
\makeatletter
\let\@tum@activate@layout@\relax
\makeatother
\small
\renewcommand{\arraystretch}{1.0}
\setlength{\tabcolsep}{4pt}
\setlength{\LTleft}{0pt}
\setlength{\LTright}{\fill}

\begin{longtable}{@{}p{0.16\linewidth}p{0.63\linewidth}p{0.15\linewidth}@{}}
\toprule
Symbol & Meaning & SI unit \\
\midrule
\endfirsthead
\toprule
Symbol & Meaning (continued) & SI unit \\
\midrule
\endhead
\bottomrule
\endlastfoot
\multicolumn{3}{@{}l}{\textbf{Latin symbols}} \\[2pt]
\(a\) & Speed of sound & \(\mathrm{m\,s^{-1}}\) \\
\(B_q\) & Dimensionless wall heat-transfer rate; negative for wall cooling & -- \\
\(c_p\) & Specific heat at constant pressure & \(\mathrm{J\,kg^{-1}\,K^{-1}}\) \\
\(C_f\) & Skin-friction coefficient & -- \\
\(C_h,\ St\) & Stanton number, with the stated reference normalization & -- \\
\(h\) & Specific enthalpy & \(\mathrm{J\,kg^{-1}}\) \\
\(k\) & Turbulent kinetic energy per unit mass & \(\mathrm{m^2\,s^{-2}}\) \\
\(M,\ M_\tau\) & Mach number and friction Mach number \(M_\tau=u_\tau/a_w\) & -- \\
\(p\) & Static pressure & \(\mathrm{Pa}\) \\
\(Pr,\ Pr_t\) & Molecular and turbulent Prandtl numbers & -- \\
\(q_w\) & Wall heat-flux; the sign convention is specified locally & \(\mathrm{W\,m^{-2}}\) \\
\(r\) & Recovery factor & -- \\
\(R\) & Specific gas constant & \(\mathrm{J\,kg^{-1}\,K^{-1}}\) \\
\(Re_x\) & Reynolds number based on streamwise distance & -- \\
\(Re_\theta\) & Momentum-thickness Reynolds number using freestream or edge viscosity & -- \\
\(Re_{\delta2}\) & Momentum-thickness Reynolds number using wall viscosity & -- \\
\(Re_\tau\) & Friction Reynolds number based on the stated boundary-layer thickness & -- \\
\(Re_{\mathrm{cell1}}\) & First-cell acoustic Reynolds number, \(\rho_w a_w y_1/\mu_w\) & -- \\
\(T\) & Temperature; \(T_w\) is the wall temperature and \(T_r\) the recovery temperature & \(\mathrm{K}\) \\
\(u,\ U\) & Streamwise or wall-tangential velocity; locally stated averaging & \(\mathrm{m\,s^{-1}}\) \\
\(u_\tau\) & Friction velocity, \(\sqrt{|\tau_w|/\rho_w}\) & \(\mathrm{m\,s^{-1}}\) \\
\(x,\ y\) & Streamwise and wall-normal coordinates & \(\mathrm{m}\) \\
\addlinespace
\multicolumn{3}{@{}l}{\textbf{Greek symbols and mesh measures}} \\[2pt]
\(\gamma\) & Ratio of specific heats & -- \\
\(\delta\) & Boundary-layer thickness, with the stated edge criterion & \(\mathrm{m}\) \\
\(\Delta x\) & Streamwise mesh spacing & \(\mathrm{m}\) \\
\(\Delta y_1\) & Full thickness of the first wall-adjacent cell & \(\mathrm{m}\) \\
\(h\) & Full cell-centre separation \(y_R-y_L\) in the internal-face stencil analysis & \(\mathrm{m}\) \\
\(\epsilon\) & Relative tolerance on the checked gradient or coefficient-interpolation error & -- \\
\(\theta\) & Momentum thickness & \(\mathrm{m}\) \\
\(\kappa,\ \lambda\) & Thermal conductivity & \(\mathrm{W\,m^{-1}\,K^{-1}}\) \\
\(\mu\) & Dynamic viscosity & \(\mathrm{Pa\,s}\) \\
\(\nu\) & Kinematic viscosity & \(\mathrm{m^2\,s^{-1}}\) \\
\(\rho\) & Density & \(\mathrm{kg\,m^{-3}}\) \\
\(\tau_w\) & Wall shear stress & \(\mathrm{Pa}\) \\
\(\omega\) & Specific turbulence-dissipation rate in the SST model & \(\mathrm{s^{-1}}\) \\
\end{longtable}

\clearpage
For the wall-normal cells used here, the first-cell-centre distance is
\(y_1=\delta_n=\Delta y_1/2\). It is distinct from the wall-function sampling
height \(y_s\), the internal-face position \(y_f\) and the cell-centre
separation \(h\). Source figures that use \(z\) as the wall-normal coordinate
are identified in their captions.

\begin{longtable}{@{}p{0.24\linewidth}p{0.72\linewidth}@{}}
\toprule
Notation & Meaning \\
\midrule
\endfirsthead
\toprule
Notation & Meaning (continued) \\
\midrule
\endhead
\bottomrule
\endlastfoot
\multicolumn{2}{@{}l}{\textbf{Subscripts}} \\[2pt]
\(w\) & Wall value \\
\(\infty\) & Freestream value \\
\(e,\ \delta\) & Boundary-layer edge value \\
\(r\) & Recovery state \\
\(\mathrm{lam},\ t,\ T,\ \mathrm{eff}\) & Laminar, turbulent (\(t\) or \(T\)) and effective quantities; \(t\) may also denotes a tangential component where stated \\
\(1,\ 2\) & Cell or sampling indices; pre-shock and post-shock states in shock relations \\
\(s,\ m\) & Wall-function sampling location and Griffin matching location, respectively \\
\(f\) & Internal face; finest-grid value in mesh comparisons \\
\(L,\ R\) & Neighbouring cell centres in the face stencil \\
\(i,\ 99\) & Inlet condition and 99\% boundary-layer edge criterion, respectively \\
\(\mathrm{ref}\) & Specified reference value \\
\(\mathrm{DNS},\ \mathrm{RANS},\ \mathrm{SA}\) & DNS, RANS or Spalart--Allmaras reference \\
\(\mathrm{WR},\ \mathrm{WF}\) & Wall-resolved or wall-function result \\
\(\mathrm{std},\ \mathrm{target}\) & Standard evaluation or prescribed target value \\
\(\mathrm{TL},\ \mathrm{VD}\) & Trettel--Larsson or van Driest transformed quantity \\
\(\mathrm{grad},\ \mathrm{interp}\) & Gradient reconstruction or coefficient interpolation \\
\(\mathrm{poly},\ \max\) & Polynomial fit or maximum value \\
\addlinespace
\multicolumn{2}{@{}l}{\textbf{Superscripts, accents and operators}} \\[2pt]
\({}^{+}\) & Wall units: lengths scaled by \(\nu_w/u_\tau\), velocities by \(u_\tau\) \\
\(y^*\) & Semi-local wall coordinate \\
\({}^{\mathrm{OF}}\) & Native OpenFOAM output, with its stated sign convention \\
\(\overline{(\cdot)}\) & Reynolds mean of a flow quantity \\
\(\widetilde{\nu}\) & Spalart--Allmaras working variable \\
\(\boldsymbol{u},\ \boldsymbol{n}\) & Bold symbols denote vectors \\
\({}',\ {}'',\ {}'''\) & Derivatives with respect to physical wall-normal distance in the stencil analysis \\
\(D_f\) & Discrete internal-face gradient operator \\
\(\nabla,\ \partial\) & Spatial gradient and partial differentiation \\
\(\partial/\partial n_{\mathrm{out}}\) & Derivative along the outward unit normal of the fluid domain \\
\(\left.(\cdot)\right|_w\) & Evaluation at the wall \\
\(|\cdot|,\ \mathcal{O}(\cdot)\) & Absolute value and asymptotic order, respectively \\
\end{longtable}
\endgroup
